\scanpage{25}
% loc 1–2, Pour récrire
Pour \unsure{récrire|réécrire} 2), on \deleted{écrit} prend $k'=A\otimes A$, et on choisit les deux hom
% loc 2, two parallel strokes
\diagram{\begin{tikzcd}[ampersand replacement=\&]
 \deleted{\unsure{g}} \& A^{\unsure{\cdot}}\arrow[r,dash,shift left=.55ex,"g"]\arrow[r,shift right=.55ex,swap,"g'"] \& A\otimes A \& \deleted{\unsure{g'}}\quad\text{donnés par}
\end{tikzcd}}
% loc 3, definitions
\[
 g(\lambda)=\lambda\otimes1,\quad g'(\lambda)=1\otimes\lambda
\]
% loc 4–7, et on écrit
et on écrit l'équation comme une égalité de deux endom. de $\deleted{\unsure{M}_{A\otimes A}}\ M_{k'}=M_{A\otimes A}$ qu'on va écrire -- puis on écrira l'égalité de leurs composés avec $M\overset{\mathrm{can}}{\longrightarrow}M_{A\otimes A}$ -- donc finalement il faudra définir deux $k$-hom $M\mathrel{\unsure{\text{---}|\to}}M_{A\otimes A}$, savoir
% loc 8, two x maps
\[
 x\longmapsto gx'_{\deleted{\unsure{k'}}},\quad x\mathrel{\unsure{\mapstochar\text{---}|\mapsto}}g'x'_{\deleted{\unsure{k}}}
 \quad\text{(où }x'_{\unsure{k|M}}\text{ est l'image de }x\text{ dans }M_{k'}=M_{A\otimes A}\text{).}
\]

% loc 9–10, Or
Or
\[
 (gg')_M x'=(gg')_M^{\unsure{0|\circ}}x=(\mathrm{id}_M\otimes\unsure{gg|gg'})\circ u_0 \quad:
\]
% loc 10, upper box and element map
\diagram{\begin{tikzcd}[ampersand replacement=\&,column sep=large,execute at end picture={\draw ([xshift=-4pt,yshift=8pt]p.north west) -- ([xshift=9pt,yshift=8pt]q.north east) -- ([xshift=9pt,yshift=-6pt]r.south east) -- ([xshift=-4pt,yshift=-6pt]p.west|-r.south) -- cycle;}]
 \& |[alias=p]|M\arrow[r,"u_0"] \& |[alias=q]|M\otimes A\arrow[d,swap,"\mathrm{id}_M\otimes\pi"]\arrow[d,phantom,shift left=5ex,"\deleted{\unsure{(u_0\otimes\mathrm{id}_A)}}"] \& \deleted{\unsure{y\otimes\lambda}}\arrow[d,maps to]\\
 (gg')_M^0\arrow[ur,shorten >=6pt] \& \& |[alias=r]|M\otimes A\otimes A \& \deleted{\unsure{y\otimes\lambda\otimes1}}
\end{tikzcd}}
% loc 11, question
\[
 g_M(g'_Mx')=\,?
\]
% loc 12–13, Calculons and interline reading
Calculons $\underset{\substack{\unsure{\shortparallel|''}\\\unsure{(g'_M)^0x^{*}|(g'_M)^0x}}}{g'_Mx'}\in M_{k'}$, c'est $(\mathrm{id}_M\otimes g')\deleted{\unsure{(}}(u_0(x))$ :
\diagram{\begin{tikzcd}[ampersand replacement=\&]
 M\arrow[r,"u_0"] \& M\otimes A\arrow[d] \& y\otimes\lambda\arrow[d,maps to]\\
 \& M\otimes A\otimes A \& y\otimes1\otimes\lambda
\end{tikzcd}}
% loc 14–15, Donc and repeated par
Donc $g_M(g'_Mx')$ s'obtient en appliquant au résultat $g_M$, donné par par
% loc 16–17, g on two simple tensors
\[
\begin{aligned}
 &g_M(y\otimes\lambda\otimes\mu)=g_M^0(y)\,\lambda\otimes\mu
 =(\mathrm{id}_M\otimes g)(u_0(y))\,\lambda\otimes\mu\\
 &\text{\unsure{On a donc|en. donc}}\qquad g_M(y\otimes1\otimes\lambda)
 =\underset{\in M\otimes A\otimes1}{\underline{(\mathrm{id}_M\otimes g)u_0(y)}}\,1\otimes\lambda 
\end{aligned}
\]
% loc 18, lower box
donc
\diagram{\begin{tikzcd}[ampersand replacement=\&,column sep=large,execute at end picture={\draw ([xshift=-8pt,yshift=4pt]a.north west) -- ([xshift=15pt,yshift=4pt]a.north east) to[out=-85,in=180] ([xshift=26pt,yshift=-5pt]a.east) -- ([xshift=5pt,yshift=-5pt]c.east|-a.east) -- ([xshift=9pt,yshift=-5pt]c.south east) -- ([xshift=-8pt,yshift=-5pt]b.south west) -- cycle;}]
 \& |[alias=a]|M\arrow[d,"u_0"] \& \\
 (g_Mg'_M)^0: \& |[alias=b]|M\otimes A\arrow[r,"u_0\otimes\mathrm{id}_A"] \& |[alias=c]|\underbrace{M\otimes A}\otimes A\\
 \& \deleted{\unsure{y\otimes\lambda}} \&
\end{tikzcd}}

% loc 19–20, Donc l'axiome
Donc l'axiome \deleted{\unsure{de|1}} 2 s'exprime comme
\diagram{\begin{tikzcd}[ampersand replacement=\&,column sep=large]
 \tikz[baseline=(t.base)]{\node[draw,circle,inner sep=1.6pt](t){$2'$};} \& M\arrow[r,"u_0"]\arrow[d,swap,"u_0"] \& M\otimes A\arrow[d,"\mathrm{id}_M\otimes\pi"] \& \text{\uline{commutatif}}\\
 \& M\otimes A\arrow[r,"u_0\otimes\mathrm{id}_A"] \& M\otimes A\otimes A \&
\end{tikzcd}}
% archive sheet number: 24
